\ifdefined\HMTIncludeParticleCore
\chapter{Secciones persistentes y propagación finita}
\label{chap:particula}

La noción matemática de partícula se introduce después de la estructura de
frontera y antes de especializar su dinámica. Su soporte primario es una
extensión superviviente del sistema de estados TPK\@. Una trayectoria
orientada es la sombra secuencial de esa extensión; sobre ella, los grados
internos forman un fibrado discreto. El propagador, la suma de caminos y las
representaciones de Clifford se incorporan después de construir la recta de
acción. Esta jerarquía impide que una partícula se defina mediante la elección
externa de una fila o de una ruta y separa el objeto matemático de su posterior
identificación con una especie física.

\section{Extensión superviviente y proyección observable}

Sea \((\mathcal S_m,\rho_{m+1,m})\) el sistema inverso de estados que
superan los cierres hasta profundidad \(m\). Un estado global es una familia
compatible
\[
 x=(x_m)_m\in\Omega_{\HMT}^{\rm surv}
 :=\varprojlim_m\mathcal S_m.
\]
Una carta secuencial de profundidad \(R\) produce una ruta
\[
 \gamma_x^{[R]}=\operatorname{sh}_R(x)
 =(c_0\longrightarrow\cdots\longrightarrow c_R).
\]
El símbolo \(\operatorname{sh}\) denota una lectura de sombra: diferentes
cartas pueden presentar un mismo estado global, mientras conserven hoja,
orientación, borde y registro. En el
\cref{chap:extension-ruta-caracter} se construye el registro enriquecido que
hace descender el carácter de acción desde estas presentaciones a la
extensión.

\section{Partícula como sección compatible de un fibrado discreto}

Sea \(\Gamma_{\rm TPK}\) el grafo cuyos vértices son estados completos del
núcleo de fase topológica y cuyas aristas son transformaciones admisibles.
Para \(x\in\Omega_{\HMT}^{\rm surv}\) y su trayectoria observada
\[
 \gamma_x=(c_0\longrightarrow c_1\longrightarrow\cdots
 \longrightarrow c_R),
\]
sea \(\operatorname{Typ}(c_t)\) el conjunto finito de tipos aritméticos,
orientaciones y memorias sobre el estado \(c_t\), y sea
\(E_{c_t}=\mathbb C[\operatorname{Typ}(c_t)]\) el espacio vectorial complejo
libre generado por esos tipos. El fibrado discreto restringido a \(\gamma_x\) es
\[
 \mathcal B_{\gamma_x}=\bigsqcup_{t=0}^{R}E_{c_t}
 \longrightarrow\{0,\ldots,R\}.
\]
Cada arista \(c_t\to c_{t+1}\) induce una aplicación lineal de transporte
\(\tau_t:E_{c_t}\to E_{c_{t+1}}\).

\begin{definition}[Partícula matemática]
Una partícula de Mecánica Dimensional es una tupla
\[
 \mathfrak p=(x,\gamma_x,\mathcal B_{\gamma_x},s,n,U),
\]
donde \(x\in\Omega_{\HMT}^{\rm surv}\) es la extensión superviviente y
\(\gamma_x=\operatorname{sh}_R(x)\) su sombra en la carta declarada; además,
\[
 s:\{0,\ldots,R\}\longrightarrow\mathcal B_{\gamma_x},\qquad
 s(t)\in E_{c_t},\qquad s(t+1)=\tau_t(s(t)),
\]
es una sección compatible;
\[
 n:\bigsqcup_{t=0}^{R}\operatorname{Modes}(E_{c_t})
 \longrightarrow\mathbb N
\]
es su función de ocupación, donde
\(\operatorname{Modes}(E_{c_t}):=\operatorname{Typ}(c_t)\) denota la base
distinguida del espacio vectorial libre; y
\[
 U_{t_2,t_1}:E_{c_{t_1}}\longrightarrow E_{c_{t_2}}
\]
es una familia de propagadores lineales que entrelaza los transportes,
\(U_{t,t}=\operatorname{id}_{E_{c_t}}\),
\(U_{t+1,t}=\tau_t\), y satisface
\[
 U_{t_3,t_2}\circ U_{t_2,t_1}=U_{t_3,t_1}
 \qquad(0\leq t_1\leq t_2\leq t_3\leq R).
\]
Cada composición está definida entre las fibras correspondientes.
La identificación de \(\mathfrak p\) con una especie
física requiere, además, una representación del grupo de simetría pertinente
y un mapa desde estados HMT hacia los estados de dicha representación.
\end{definition}

La extensión conserva la genealogía y el registro; la ruta expone su lectura
ordenada; los grados internos pertenecen a las fibras, la ocupación a \(n\) y
la evolución a \(U\). La partícula es la compatibilidad de estos niveles, no
uno de ellos por separado. En particular, la ruta no es escogida por la
partícula: es una proyección de la extensión que la realiza.
\fi

\ifdefined\HMTIncludeParticleQuantum
\section{Expansión matricial como suma finita de caminos}
\index{camino!expansión matricial finita}

\begin{theorem}[Expansión matricial como suma finita de caminos]
Sea \(U\) una matriz indexada por un conjunto finito de estados. Para
\(R\geq1\),
\[
 (U^R)_{fi}
 =\sum_{\substack{\gamma:i\to f\\|\gamma|=R}}
 \prod_{\ell\in\gamma}U_\ell.
\]
\end{theorem}
\begin{proof}
La multiplicación matricial proporciona
\[
 (U^R)_{fi}=\sum_{x_1,\ldots,x_{R-1}}
 U_{fx_{R-1}}\cdots U_{x_2x_1}U_{x_1i}.
\]
Cada tupla de índices intermedios determina una trayectoria de longitud \(R\),
y cada trayectoria aparece una vez. La igualdad es enteramente finita; no
presupone un límite funcional ni una integral de caminos continua
\cite{Feynman1948}.
\end{proof}

La identidad anterior es una expansión matricial finita exacta: compone los
pesos ordenados de historias HMT ya construidas. Cuando \(U\) realiza un
operador de evolución, esta expansión proporciona el antecedente discreto de
la suma de caminos de Feynman. El paso desde la suma finita hasta una integral
sobre genealogías y, después, hasta un propagador físico se formula por
separado en la \cref{sec:integral-genealogica-caminos-rev2}.

\begin{corollary}[Amplitud global de la banda finita]
\label{cor:amplitud-global-banda-finita}
Sea \(X\) el conjunto finito de estados, sea
\(\psi_0\in\mathbb C^X\) y defínase \(\psi_R=U^R\psi_0\). Entonces, para
cada estado final \(f\),
\[
 \psi_R(f)
 =\sum_{i\in X}
   \sum_{\substack{\gamma:i\to f\\|\gamma|=R}}
   w_U(\gamma)\,\psi_0(i),
 \qquad
 w_U(\gamma):=
 \prod_{t=1}^{R}U_{x_t x_{t-1}}.
\]
Si \(U\) es unitario, además,
\[
 \|\psi_R\|_2=\|\psi_0\|_2.
\]
\end{corollary}

\begin{proof}
La primera identidad resulta de aplicar el teorema anterior a cada
coeficiente de
\[
 (U^R\psi_0)(f)=\sum_{i\in X}(U^R)_{fi}\psi_0(i).
\]
Si \(U\) es unitario, también lo es \(U^R\), y la segunda igualdad se sigue
de la invariancia de la norma hermítica.
\end{proof}

El vector \(\psi_R=U^R\psi_0\) es la función de onda HMT a profundidad finita
\(R\) sobre el espacio discreto de estados declarado. Cada componente
\(\psi_R(f)\) suma coherentemente las amplitudes de todas las historias que
terminan en \(f\). Si \(U\) es unitario, la norma total se conserva. La
realización como función de onda de una especie física concreta consiste
después en especificar la representación que identifica ese espacio HMT, su
propagador y sus observables con el sector físico correspondiente.

\section{Representación real de Clifford y realificación de Dirac}
\label{sec:dirac-real-hmt-md}
\index{Dirac!realificación}\index{Clifford!representación real}
\index{función de onda!amplitud global finita}

Conviene separar dos resultados algebraicos. El primero exhibe una
representación real de la relación de Clifford lorentziana. El segundo
realifica cualquier ecuación compleja de Dirac que haya sido fijada de
antemano. Ninguno selecciona por sí solo la derivada covariante, la conexión,
el espacio físico de espinores ni la masa.

Sean
\[
 I_2=\begin{pmatrix}1&0\\0&1\end{pmatrix},
 \qquad
 J_2=\begin{pmatrix}0&1\\-1&0\end{pmatrix},
 \qquad
 R_2=\begin{pmatrix}0&1\\1&0\end{pmatrix},
 \qquad
 Z_2=\begin{pmatrix}1&0\\0&-1\end{pmatrix}.
\]
Estas matrices satisfacen
\[
 J_2^2=-I_2,\qquad R_2^2=Z_2^2=I_2,
\]
y anticonmutan dos a dos.

\begin{theorem}[Representación real explícita de
\(\mathrm{Cl}(3,1)\)]
\label{thm:clifford-real-31}
Las matrices reales
\[
 \gamma_0=J_2\otimes I_2,\qquad
 \gamma_1=R_2\otimes I_2,\qquad
 \gamma_2=Z_2\otimes R_2,\qquad
 \gamma_3=Z_2\otimes Z_2
\]
verifican
\[
 \gamma_\mu\gamma_\nu+\gamma_\nu\gamma_\mu
 =2\eta_{\mu\nu}I_4,
 \qquad
 \eta=\operatorname{diag}(-1,1,1,1).
\]
Por consiguiente, determinan una representación real de
\(\mathrm{Cl}(3,1)\cong M_4(\mathbb R)\).
\end{theorem}

\begin{proof}
Los cuadrados de los productos de Kronecker dan
\(\gamma_0^2=-I_4\) y \(\gamma_k^2=I_4\) para \(k=1,2,3\).
Para índices distintos, la anticomutación de \(J_2,R_2,Z_2\), o la de
los segundos factores \(R_2,Z_2\) cuando los primeros coinciden, anula el
anticomutador. Éstas son exactamente las relaciones universales que definen
el álgebra de Clifford de signatura \((-+++)\).
\end{proof}

Para enunciar el segundo resultado sin identificar objetos de dominios
distintos, sea \(A=A_1+\mathrm iA_2\in M_n(\mathbb C)\), con
\(A_1,A_2\in M_n(\mathbb R)\), y defínase
\[
 \rho_{\rm sp}(A):=
 \begin{pmatrix}
  A_1&-A_2\\
  A_2&A_1
 \end{pmatrix},
 \qquad
 J_{\rm sp}:=\rho_{\rm sp}(\mathrm iI_n)
 =\begin{pmatrix}
  0&-I_n\\
  I_n&0
 \end{pmatrix}.
\]

\begin{theorem}[Realificación exacta de una ecuación de Dirac fijada]
\label{thm:realificacion-dirac}
La aplicación \(\rho_{\rm sp}\) es un homomorfismo real inyectivo y satisface
\[
 \rho_{\rm sp}(AB)=\rho_{\rm sp}(A)\rho_{\rm sp}(B),
 \qquad
 \rho_{\rm sp}(A^\dagger)=\rho_{\rm sp}(A)^{\mathsf T},
 \]
\[
 J_{\rm sp}^{\,2}=-I_{2n},
 \qquad
 [J_{\rm sp},\rho_{\rm sp}(A)]=0.
\]
Sean \(\gamma^\mu\in M_n(\mathbb C)\) una representación compleja de
Clifford, \(D_\mu\) operadores complejos lineales ya fijados y
\(\Psi=(\operatorname{Re}\psi,\operatorname{Im}\psi)^{\mathsf T}\) la
realificación de \(\psi\). Supóngase que cada \(D_\mu\) es complejo lineal y
compatible con esta realificación. Entonces
\[
 \bigl(\mathrm i\gamma^\mu D_\mu-mI_n\bigr)\psi=0
\]
es equivalente a
\[
 \rho_{\rm sp}\!\left[
  \bigl(\mathrm i\gamma^\mu D_\mu-mI_n\bigr)\psi
 \right]=0
 \quad\Longleftrightarrow\quad
 \bigl(
 J_{\rm sp}\rho_{\rm sp}(\gamma^\mu D_\mu)
 -mI_{2n}
 \bigr)\Psi=0,
\]
y, por multiplicatividad,
\[
 \bigl(
 J_{\rm sp}\Gamma^\mu D_\mu^{\mathbb R}
 -mI_{2n}
 \bigr)\Psi=0,
 \qquad
 \Gamma^\mu:=\rho_{\rm sp}(\gamma^\mu),
\]
donde \(D_\mu^{\mathbb R}\) es la realificación de \(D_\mu\).
Para \(n=4\), esta segunda construcción actúa en dimensión real ocho; no es
la representación real irreducible de dimensión cuatro del teorema anterior.
\end{theorem}

\begin{proof}
La multiplicación por bloques demuestra directamente que
\(\rho_{\rm sp}\) preserva sumas, productos, adjuntos y la identidad, y que
su núcleo es nulo. La matriz \(J_{\rm sp}\) representa la multiplicación compleja por
\(\mathrm i\); por ello su cuadrado es \(-I_{2n}\) y conmuta con todo
operador complejo lineal realificado. Aplicar \(\rho_{\rm sp}\) al operador
\(\mathrm i\gamma^\mu D_\mu-mI_n\) produce exactamente el operador real del
enunciado, de modo que ambos núcleos se corresponden bajo
\(\psi\mapsto\Psi\).
\end{proof}

La raíz ternaria
\(J_W\in\operatorname{End}_{\mathbb F_3}(\mathbb F_3^6)\) construida en
\cref{chap:algebras-cuadraticas-orden-nueve} y la estructura
\(J_{\rm sp}\) sobre un espacio espinorial real son objetos distintos.
Comparten la relación polinómica \(X^2+1=0\), pero no se identifican sin un
entrelazador que cambie correctamente de dominio y de escalares.

La expansión finita de caminos, el
\cref{cor:amplitud-global-banda-finita,thm:clifford-real-31,%
thm:realificacion-dirac} son resultados algebraicos exactos. Juntos
distinguen expansión de caminos, representación real de
\(\mathrm{Cl}(3,1)\) y realificación de una ecuación compleja de Dirac: tres
objetos coordinados, pero no intercambiables.
\fi
